% Source: https://artofproblemsolving.com/community/c3h1382206_pentagon_angle_problem
% Stable identity: aops-topic-1382206-q01
The interior angles of a pentagon add to $540^\circ$, so each of $\angle W,\angle X,\angle Y,\angle Z$ is $\dfrac{540^\circ-100^\circ}{4}=110^\circ$.

\begin{asy}
usepackage("amsmath");
size(7cm);
defaultpen(fontsize(10pt));
pair V=(0,0),Wp=(1,0),X=(1.34202,0.93969),Y=(0.57598,1.58248),Z=(-0.20003,1.13445),Q=(-0.88245,0.74046);
draw(V--Wp--X--Y--Z--cycle,linewidth(0.9));
draw(Z--Q,blue); draw(Q--V,blue,MidArrow(5)); draw(Y--X,MidArrow(5));
dot(Q);
label("$V$",V,SW); label("$W$",Wp,SE); label("$X$",X,E); label("$Y$",Y,N); label("$Z$",Z,N); label("$Q$",Q,W);
draw(arc(V,0.25,degrees(Z-V),degrees(Q-V)),red); label("$40^\circ$",V+0.42*dir(115),fontsize(8pt)+red);
draw(arc(Z,0.13,degrees(Q-Z),degrees(V-Z)),red); label("$70^\circ$",Z+0.29*dir(245),fontsize(8pt)+red);
draw(arc(Q,0.13,degrees(V-Q)-360,degrees(Z-Q)),red); label("$70^\circ$",Q+0.3*dir(-6),fontsize(8pt)+red);
label("\mbox{}",(min(currentpicture,user=true)+max(currentpicture,user=true))/2,black); // keeps the text after the figure black when LaTeX typesets the labels inline
\end{asy}

In triangle $QVZ$:
\begin{itemize}
\item $\angle QZV=180^\circ-\angle YZV=180^\circ-110^\circ=70^\circ$, since $Q$, $Z$, $Y$ are collinear.
\item $QV\parallel YX$ and line $QY$ is a transversal, with $V$ and $X$ on the same side of it, so the co-interior angles add to $180^\circ$: $\angle VQY+\angle QYX=180^\circ$. Since $\angle QYX=\angle ZYX=110^\circ$, we get $\angle VQZ=70^\circ$.
\end{itemize}
Therefore $m\angle QVZ=180^\circ-70^\circ-70^\circ=\boxed{40^\circ}$.
